\section{The ConvergenceSweep Alternative}
\label{sec:sweep}

\subsection{Algorithm and Stopping Criterion}

The \cs{} algorithm~\citep{b16paper} is a deterministic forward sweep
over seeds: starting at $s_0 = \mathrm{SHA\text{-}256}(\texttt{census\_release\_id}
\,\|\, \texttt{"DIA\_SEED\_V1"}) \bmod 2^{31}$,
it evaluates each seed in order and halts when $T = 600$ consecutive
seeds produce no improvement in normalised edge cut
$\mathrm{EC}_{\rm norm} = \mathrm{EC} / \sqrt{\min(i, k-i)}$.

The runtime is:
\[
  O(T \cdot k \cdot m_{\rm METIS}(n, k)),
\]
where $m_{\rm METIS}(n,k) = O(n \log n)$ is the METIS runtime for one
partitioning call.
For $T = 600$, $k \leq 52$, $n \leq 8{,}000$, this is under $10^8$
operations---well within 60 seconds on a 2024-generation machine.

\subsection{ConvergenceSweep vs.\ MCMC: Conceptual Comparison}

\begin{table}[h]
\centering
\caption{Conceptual comparison of \cs{} and certified \recom{} ensemble.}
\label{tab:comparison}
\begin{tabular}{lcc}
\toprule
Property & \cs{} & Certified \recom{} \\
\midrule
Output & One optimal plan & Distribution over plans \\
Determinism & Yes & Only with fixed seed \\
Certified stopping & Yes ($T=600$) & Yes (G.4 diagnostics) \\
Runtime (TX) & $<$60 sec & $\approx$100 min \\
Runtime (CA) & $<$60 sec & $\approx$300 min \\
Theory basis & Gumbel tail model & Spectral gap \\
Task & Plan generation & Plan evaluation \\
\bottomrule
\end{tabular}
\end{table}

\subsection{Complexity Comparison}

Let us compare the asymptotic complexity of the two approaches for
producing a certifiably good redistricting plan:

\paragraph{\cs{}.}
$O(T \cdot k \cdot n \log n)$ with $T = 600$ (constant).
For the hardest state (CA, $k = 52$, $n = 8{,}057$):
$600 \times 52 \times 8{,}057 \times \log(8{,}057) \approx 1.8 \times 10^9$ operations.
Wall time: $<$60 seconds.

\paragraph{Certified \recom{} for plan generation.}
If one wanted to use \recom{} to generate the ``most typical'' plan
(the mode of the distribution), one would need to:
(1) run the chain for $\hat{t}_{\rm mix}$ steps to reach stationarity,
(2) draw $\ess$ independent samples,
(3) select the sample with the minimum edge cut.
For CA: $50{,}000 \times (8{,}057 \times \log(8{,}057)) \approx 6 \times 10^9$ operations
just for mixing, plus sampling overhead.
Wall time: $\approx$300 minutes.

\paragraph{Conclusion.}
For plan generation, \cs{} is $O(T/(n_{\rm min} \cdot k))$ faster than
certified \recom{}---approximately $300\times$ faster for California.
The stopping certificate ($T = 600$ consecutive non-improving seeds) is
more transparent than an MCMC convergence argument.

\subsection{Epistemic Comparison}

The Gumbel tail model underlying $T = 600$ makes an explicit probabilistic
claim: $P(\text{tail} > 600) < 0.001$ for all states (B.16 Section 3).
This claim is falsifiable: given the 50-state sweep data from B.7, the
Gumbel fit can be checked against the empirical distribution of tails.

The MCMC mixing claim makes a weaker statement: $\rhat < 1.1$ at step
$\hat{t}_{\rm mix}$.
This is a necessary condition for approximate convergence but not a
sufficient condition for total-variation convergence.
The spectral gap bound $\Omega(1/n^2)$ gives the theoretical worst case,
which is far from the empirical mixing time.

Neither approach is ``more rigorous'' in an absolute sense; they formalize
different aspects of the stopping-criterion problem.
\cs{} is appropriate for deterministic plan generation; ensemble sampling
with G.4 diagnostics is appropriate for distributional inference.
